\subsection{Subsheaves of a sheaf}

Let B be a topological space. Let \(\mathcal{F} = (\mathcal{F}(U), f_{UV})\) be a presheaf on B, relative to a base \(\mathcal{B}\) of the topology of B.

Suppose that, for every open set \(U \in \mathcal{B}\), there is a subset \(\mathcal{L}(U)\) of \(\mathcal{F}(U)\) . If \(f_{\mathrm{UV}}(\mathcal{L}(V)) \subset \mathcal{L}(U)\) for every pair \((U, V)\) of elements of \(\mathcal{B}\) such that \(U \subset V\), the pair \(\mathcal{L} = ((\mathcal{L}(U))_{U \in \mathcal{B}}, (f'_{UV}))\) , where \(f'_{UV} \colon \mathcal{L}(V) \to \mathcal{L}(U)\) is the map deduced from \(f_{UV}\) , is a presheaf. Such a presheaf is called a subpresheaf of \(\mathcal{F}\). Since the maps \(f'_{UV}\) are determined by the data of the presheaf F and the family \((\mathcal{L}(U))_{U \in \mathcal{B}}\) , one also says, by abuse of language that the family \((\mathcal{L}(U))_{U \in \mathcal{B}}\) is a subpresheaf of F.

Suppose now that \(\mathcal{F}\) is a sheaf on B, and, for every open subset U of B, let \(\mathcal{L}(\mathrm{U})\) be a subset of \(\mathcal{F}(\mathrm{U})\). For \((\mathcal{L}(\mathrm{U}))_{\mathrm{U}\in \mathcal{B}}\) to be a subpresheaf of \(\mathcal{F}\), and for this presheaf to be a sheaf, it is necessary and sufficient that the following condition be satisfied:

\begin{enumerate}
\def\labelenumi{(\Alph{enumi})}
\setcounter{enumi}{5}
\tightlist
\item
  Let \((\mathrm{U}_{i})_{i\in\mathrm{I}}\) be a family of open sets of B, let U be its union, and let s be an element of \(\mathcal{F}(\mathrm{U})\). For s to belong to \(\mathcal{L}(\mathrm{U})\) , it is necessary and sufficient that:
\end{enumerate}

\[
\text{for every } i \in \mathrm{I},\quad f_{\mathrm{U}_{i}\mathrm{U}}(s)\in\mathcal{L}(\mathrm{U}_{i}).
\]

Indeed, if condition (F) holds, we have \(f_{\mathrm{UV}}(\mathcal{L}(\mathrm{V})) \subset \mathcal{L}(\mathrm{U})\) for every pair (U, V) of open sets of B such that U ⊂ V and the properties \(F_{1}\) and \(F_{2}\) relative to the subpresheaf \(\mathcal{L}(\mathrm{U})\) follow from the analogous properties relative to the sheaf F. The converse is immediate.

When condition (F) is satisfied, one says that \((\mathcal{L}(\mathrm{U}))\) is a subsheaf of the sheaf F.

